\documentclass{ximera}
%herschreven 23 september 2026
\title{Theorie: Elektrisch veld}
\author{Daan Lipkens}

\begin{document}
\begin{abstract}
    %Introductie van het veldbegrip in de elektrostatica.
\end{abstract}

\maketitle
\label{hoofdstuk:elektrisch_veld}
% ----------------------------------------------------------------------
%          inleiding hoofdstuk
% ----------------------------------------------------------------------
\section*{Inleiding}
In het vorige hoofdstuk bestudeerden we de coulombkracht: een kracht die op afstand inwerkt, zonder dat er fysiek contact nodig is tussen twee geladen voorwerpen. Deze \emph{actie op afstand} riep lange tijd vragen op, tot de Britse wetenschapper Michael Faraday (figuur~\ref{fig:faraday_portret}) het concept van een \textbf{veld} introduceerde.

%---------------------------
% AFBEELDING: portret Michael Faraday
% bron: https://commons.wikimedia.org/wiki/File:M_Faraday_Th_Phillips_oil_1842.jpg
\begin{figure}[ht!]
    \captionsetup{width=0.85\textwidth}
    \begin{center}
        \includegraphics[width=0.3\textwidth]{pic5/elektriciteit/veld/faraday_portret.jpg}
        \fotobron{Thomas Phillips (1842)}{publiek domein}{https://commons.wikimedia.org/wiki/File:M_Faraday_Th_Phillips_oil_1842.jpg}
    \end{center}
    \caption{Michael Faraday (1791--1867), geschilderd door Thomas Phillips.}
    \label{fig:faraday_portret}
\end{figure}
%---------------------------


Volgens dit model creëert elke lading in de ruimte rondom zich een elektrisch veld. Een andere lading die zich in dat veld bevindt, ondervindt daardoor een kracht. Op die manier verloopt de krachtwerking niet langer rechtstreeks van lading tot lading, maar via het veld dat aanwezig is in de tussenliggende ruimte.

In dit hoofdstuk leer je hoe je het elektrisch veld van één of meerdere puntladingen berekent, hoe je een veld grafisch voorstelt met behulp van veldlijnen, en hoe geleiders en isolatoren zich gedragen wanneer ze zich in een elektrisch veld bevinden.

\end{document}